
\chapter{Contexte Général du Projet}
\label{chap:contexte}

\section{Introduction}

Ce chapitre présente le cadre général dans lequel s'inscrit notre projet de fin d'études. Nous commençons par définir le contexte et la problématique, puis nous analysons l'existant avant de proposer notre solution. Enfin, nous détaillons la méthodologie de travail adoptée ainsi que le planning du projet.

\section{Présentation de l'Organisme}

Ce projet de fin d'études a été réalisé dans le cadre d'un projet académique personnel visant à développer une plateforme innovante exploitant les dernières avancées en intelligence artificielle. L'objectif est de créer un outil qui répond aux besoins réels des étudiants, professionnels et chercheurs dans leur quête de connaissances et d'apprentissage.

\section{Cadre du Projet}

\subsection{Contexte Général}

L'intelligence artificielle connaît une évolution fulgurante, particulièrement dans le domaine du traitement du langage naturel. Les modèles de langage de grande taille (LLM) comme GPT-4, Claude et Gemini ont démontré des capacités impressionnantes dans la compréhension et la génération de texte, ouvrant ainsi de nouvelles opportunités pour l'éducation et l'apprentissage.

Le secteur de l'éducation en ligne se transforme rapidement avec l'intégration de ces technologies. Les apprenants recherchent des outils plus intelligents, plus interactifs et plus personnalisés pour accompagner leur parcours éducatif.

\subsection{Motivation}

Les principales motivations de ce projet sont :

\begin{itemize}
    \item Démocratiser l'accès aux modèles d'IA les plus performants via une interface unifiée
    \item Faciliter l'apprentissage en intégrant des fonctionnalités complémentaires (analyse de documents, visualisations, recommandations)
    \item Créer un environnement collaboratif favorisant le partage de connaissances
    \item Proposer des outils de génération automatique de contenu éducatif
\end{itemize}

\section{Problématique}

\subsection{Constat}

Actuellement, les utilisateurs souhaitant exploiter la puissance des modèles de langage font face à plusieurs défis :

\begin{enumerate}
    \item \textbf{Fragmentation des plateformes} : Pour accéder à GPT-4, Claude et Gemini, il faut créer des comptes sur trois plateformes différentes (OpenAI, Anthropic, Google)
    \item \textbf{Coûts d'abonnement multiples} : Chaque plateforme nécessite un abonnement séparé, ce qui peut représenter un coût élevé
    \item \textbf{Manque d'outils complémentaires} : Les plateformes existantes se limitent souvent au chat textuel sans proposer d'analyse de documents ou de génération de visualisations
    \item \textbf{Absence de dimension collaborative} : Les utilisateurs ne peuvent pas partager facilement leurs découvertes ou interagir avec une communauté
    \item \textbf{Ressources éducatives dispersées} : Il n'existe pas de lien direct entre les conversations avec l'IA et les ressources vidéo pertinentes
\end{enumerate}

\subsection{Questions de Recherche}

Notre projet cherche à répondre aux questions suivantes :

\begin{itemize}
    \item Comment créer une interface unifiée permettant d'exploiter plusieurs modèles LLM simultanément ?
    \item Comment intégrer intelligemment l'analyse de documents dans un workflow d'apprentissage ?
    \item Comment générer automatiquement des visualisations pertinentes à partir de données ?
    \item Comment recommander des ressources vidéo contextuellement pertinentes ?
    \item Comment favoriser les échanges et le partage de connaissances au sein d'une communauté d'utilisateurs ?
\end{itemize}

\section{Étude et Critique de l'Existant}

\subsection{Étude de l'Existant}

Plusieurs plateformes proposent des fonctionnalités similaires à celles envisagées pour SynapseAI :

\begin{table}[H]
\centering
\caption{Comparaison des plateformes existantes}
\label{tab:existant}
\begin{tabular}{|l|c|c|c|c|}
\hline
\textbf{Plateforme} & \textbf{Multi-LLM} & \textbf{Analyse Docs} & \textbf{Communauté} & \textbf{Graphiques} \\
\hline
ChatGPT (OpenAI) & Non & Oui & Non & Non \\
\hline
Claude (Anthropic) & Non & Oui & Non & Non \\
\hline
Bard/Gemini (Google) & Non & Oui & Non & Non \\
\hline
Poe & Oui & Limité & Non & Non \\
\hline
Notion AI & Non & Oui & Oui & Non \\
\hline
\textbf{SynapseAI} & \textbf{Oui} & \textbf{Oui} & \textbf{Oui} & \textbf{Oui} \\
\hline
\end{tabular}
\end{table}

% Placeholder pour figure comparative
\begin{figure}[H]
    \centering
    \fbox{\parbox{0.8\textwidth}{\centering [PLACEHOLDER - Figure comparative des plateformes existantes]\\Insérer un graphique comparatif montrant les fonctionnalités de chaque plateforme}}
    \caption{Comparaison visuelle des fonctionnalités des plateformes existantes}
    \label{fig:comparaison-plateformes}
\end{figure}

\subsubsection{ChatGPT (OpenAI)}

ChatGPT est l'une des plateformes les plus populaires exploitant le modèle GPT-4. Elle offre :
\begin{itemize}
    \item Une interface de chat conversationnelle
    \item L'analyse de documents (PDF, images)
    \item La génération de code
    \item Un système de plugins (version payante)
\end{itemize}

\textbf{Limitations} : Accès limité à un seul modèle, absence de fonctionnalités communautaires, pas de recommandations de ressources externes.

\subsubsection{Claude (Anthropic)}

Claude propose une interface similaire avec des capacités d'analyse de documents avancées :
\begin{itemize}
    \item Traitement de longs documents (jusqu'à 100k tokens)
    \item Analyse précise et structurée
    \item Conversations contextuelles
\end{itemize}

\textbf{Limitations} : Un seul modèle disponible, pas d'intégration avec d'autres services, interface minimaliste.

\subsubsection{Poe}

Poe de Quora offre un accès à plusieurs modèles d'IA :
\begin{itemize}
    \item Accès à GPT-4, Claude, Gemini et d'autres modèles
    \item Création de bots personnalisés
    \item Interface unifiée
\end{itemize}

\textbf{Limitations} : Fonctionnalités avancées limitées, pas de génération de graphiques, absence de communauté, interface peu intuitive.

\subsection{Solution Proposée}

Face aux limitations des solutions existantes, nous proposons \textbf{SynapseAI}, une plateforme complète qui intègre :

\begin{enumerate}
    \item \textbf{Interface Multi-LLM Unifiée}
    \begin{itemize}
        \item Accès à GPT-4, Claude et Gemini depuis une seule interface
        \item Possibilité de comparer les réponses de différents modèles
        \item Sélection du modèle en fonction du type de requête
    \end{itemize}
    
    \item \textbf{Analyse Intelligente de Documents}
    \begin{itemize}
        \item Upload de documents PDF, DOCX, TXT
        \item Extraction automatique des concepts clés
        \item Génération de résumés structurés
        \item Interaction conversationnelle avec le contenu
    \end{itemize}
    
    \item \textbf{Génération de Visualisations}
    \begin{itemize}
        \item Création automatique de graphiques (lignes, barres, circulaires)
        \item Basé sur Chart.js pour des visualisations interactives
        \item Export et personnalisation des graphiques
    \end{itemize}
    
    \item \textbf{Recommandations YouTube Contextuelles}
    \begin{itemize}
        \item Analyse des sujets de conversation
        \item Extraction de mots-clés pertinents
        \item Recommandations de vidéos éducatives via YouTube API
        \item Intégration directe dans l'interface
    \end{itemize}
    
    \item \textbf{Espace Communautaire}
    \begin{itemize}
        \item Publication de posts et d'articles
        \item Système de commentaires
        \item Likes et interactions sociales
        \item Partage de connaissances
    \end{itemize}
    
    \item \textbf{Générateur d'Exercices}
    \begin{itemize}
        \item Création automatique d'exercices personnalisés
        \item Correction automatique
        \item Suivi de la progression
    \end{itemize}
    
    \item \textbf{Tableau de Bord Administratif}
    \begin{itemize}
        \item Gestion des utilisateurs
        \item Modération du contenu
        \item Statistiques d'utilisation
        \item Système de tickets
    \end{itemize}
\end{enumerate}

% Placeholder pour architecture globale
\begin{figure}[H]
    \centering
    \fbox{\parbox{0.9\textwidth}{\centering [PLACEHOLDER - Architecture globale de SynapseAI]\\Insérer un diagramme montrant l'architecture générale avec les différents modules}}
    \caption{Architecture globale de la plateforme SynapseAI}
    \label{fig:architecture-globale}
\end{figure}

\section{Méthodologie de Travail}

\subsection{Étude Comparative des Méthodologies}

Avant de choisir notre méthodologie de développement, nous avons comparé plusieurs approches :

\begin{table}[H]
\centering
\caption{Comparaison des méthodologies de développement}
\label{tab:methodologies}
\begin{tabular}{|l|p{4cm}|p{4cm}|p{3cm}|}
\hline
\textbf{Méthodologie} & \textbf{Avantages} & \textbf{Inconvénients} & \textbf{Adaptabilité} \\
\hline
Cascade & Structure claire, documentation complète & Rigide, difficile de revenir en arrière & Faible \\
\hline
V-Model & Tests intégrés, traçabilité & Lourd pour petits projets & Moyenne \\
\hline
\textbf{SCRUM} & Flexible, itératif, feedback rapide & Nécessite discipline & Élevée \\
\hline
Kanban & Visualisation claire, flux continu & Moins structuré & Élevée \\
\hline
XP & Qualité du code, tests & Demande expertise & Moyenne \\
\hline
\end{tabular}
\end{table}

\subsection{Méthodologie Adoptée : SCRUM}

Nous avons choisi la méthodologie agile \textbf{SCRUM} pour les raisons suivantes :

\subsubsection{Pourquoi SCRUM ?}

\begin{enumerate}
    \item \textbf{Flexibilité} : SCRUM permet d'adapter le projet aux changements et aux nouvelles exigences
    \item \textbf{Livraisons Incrémentales} : Chaque sprint produit des fonctionnalités utilisables
    \item \textbf{Feedback Rapide} : Les itérations courtes permettent d'ajuster rapidement le développement
    \item \textbf{Gestion des Risques} : Les problèmes sont identifiés et résolus rapidement
    \item \textbf{Transparence} : Suivi clair de l'avancement du projet
\end{enumerate}

% Placeholder pour cycle SCRUM
\begin{figure}[H]
    \centering
    \fbox{\parbox{0.8\textwidth}{\centering [PLACEHOLDER - Cycle SCRUM]\\Insérer un diagramme du cycle SCRUM avec Sprint Planning, Daily Scrum, Sprint Review, Sprint Retrospective}}
    \caption{Le cycle SCRUM appliqué au projet SynapseAI}
    \label{fig:cycle-scrum}
\end{figure}

\subsubsection{Les Rôles}

Dans notre projet, les rôles SCRUM sont définis comme suit :

\begin{itemize}
    \item \textbf{Product Owner} : [Nom] - Définit les besoins et priorise les fonctionnalités
    \item \textbf{Scrum Master} : [Nom] - Facilite le processus SCRUM et résout les obstacles
    \item \textbf{Development Team} : [Nom(s)] - Équipe de développement technique
    \item \textbf{Encadrant Académique} : [Nom] - Supervise et valide les livrables
\end{itemize}

\subsubsection{Organisation des Sprints}

Le projet a été décomposé en trois sprints de 4 semaines chacun :

\begin{enumerate}
    \item \textbf{Sprint 1} (Semaines 1-4) : Infrastructure et Fonctionnalités de Base
    \begin{itemize}
        \item Mise en place de l'architecture technique
        \item Système d'authentification (JWT, OAuth Google)
        \item Intégration des trois modèles LLM
        \item Interface de chat de base
        \item Gestion de l'historique des conversations
    \end{itemize}
    
    \item \textbf{Sprint 2} (Semaines 5-8) : Fonctionnalités Avancées d'IA
    \begin{itemize}
        \item Système d'analyse de documents
        \item Générateur de graphiques avec Chart.js
        \item Intégration YouTube API pour recommandations
        \item Interface Workspace pour documents
        \item Optimisation des performances
    \end{itemize}
    
    \item \textbf{Sprint 3} (Semaines 9-12) : Communauté et Administration
    \begin{itemize}
        \item Espace communautaire (Posts, Comments, Likes)
        \item Générateur d'exercices personnalisés
        \item Tableau de bord administratif
        \item Système de tickets
        \item Tests et déploiement
    \end{itemize}
\end{enumerate}

\section{Planning du Projet}

\subsection{Diagramme de Gantt}

Le planning global du projet est représenté dans le diagramme de Gantt suivant :

% Placeholder pour Gantt
\begin{figure}[H]
    \centering
    \fbox{\parbox{0.95\textwidth}{\centering [PLACEHOLDER - Diagramme de Gantt]\\Insérer un diagramme de Gantt montrant les 3 sprints avec leurs tâches principales et durées}}
    \caption{Diagramme de Gantt du projet SynapseAI}
    \label{fig:gantt}
\end{figure}

\subsection{Jalons Principaux}

\begin{table}[H]
\centering
\caption{Jalons principaux du projet}
\label{tab:jalons}
\begin{tabular}{|l|l|p{6cm}|}
\hline
\textbf{Jalon} & \textbf{Date} & \textbf{Livrables} \\
\hline
Fin Sprint 1 & Semaine 4 & Architecture + Auth + Chat Multi-LLM \\
\hline
Fin Sprint 2 & Semaine 8 & Analyse docs + Graphiques + YouTube \\
\hline
Fin Sprint 3 & Semaine 12 & Communauté + Exercices + Admin Dashboard \\
\hline
Soutenance & Semaine 14 & Rapport + Présentation + Démonstration \\
\hline
\end{tabular}
\end{table}

\section{Conclusion}

Dans ce chapitre, nous avons présenté le contexte général de notre projet SynapseAI. Nous avons identifié la problématique liée à la fragmentation des plateformes d'IA et au manque d'outils intégrés pour l'apprentissage. Notre solution propose une plateforme unifiée combinant l'accès à plusieurs modèles LLM avec des fonctionnalités avancées d'analyse, de visualisation et de collaboration.

L'adoption de la méthodologie SCRUM nous permet d'aborder le développement de manière agile et itérative, garantissant des livraisons régulières et la possibilité d'ajuster le projet selon les besoins.

Dans le chapitre suivant, nous détaillerons l'état de l'art des technologies utilisées ainsi que la spécification complète des besoins fonctionnels et non fonctionnels de la plateforme.

